% Einheit 3 von 5 · Summe mal Summe (binomische-formeln.md, Einheit 1) – Hauptnummern 20–27
\setcounter{aufgabe}{19}
\zweigkopf{Einheit 3 von 5 · Summe mal Summe}{Hier lernst du, Terme mit zwei Variablen zusammenzufassen und zwei Klammern miteinander malzunehmen · neu oder schon bekannt – je nach Buch · keine P10-Aufgabe · baut auf: Klammern auflösen (Einheit 1), Terme zusammenfassen, Terme malnehmen}{e3}

\begin{aufgabe}{Ich erkenne, welches Glied mit welchem malgenommen wird. Zeichne von jedem Glied der ersten Klammer einen Bogen zu jedem Glied der zweiten Klammer. Rechne nichts aus.}
\begin{teilezwei}
\tz \raisebox{0pt}[11mm][0pt]{\large $(x + 1) \cdot (x + 9)$} & \tz \raisebox{0pt}[11mm][0pt]{\large $(m + 3) \cdot (n + 4)$} \\
\tz \raisebox{0pt}[11mm][0pt]{\large $(x - 5) \cdot (x + 1)$} & \tz \raisebox{0pt}[11mm][0pt]{\large $(2y + 1) \cdot (y - 7)$} \\
\end{teilezwei}
\end{aufgabe}

\begin{aufgabe}{Ich kann den Wert eines Terms mit zwei Variablen berechnen. Setze $m = 4$ und $n = 1$ ein und berechne.}
\beispiel{m = 2,\ n = 5\text{:}\quad 3m + n &= 3 \cdot 2 + 5 = 11}
\begin{teilezwei}
\tz $m + n =$ \leerfeld & \tz $2m - n =$ \leerfeld \\
\tz $m \cdot n =$ \leerfeld & \tz $3m + 5n =$ \leerfeld \\
\end{teilezwei}
Setze jetzt $m = -3$ und $n = 2$ ein.
\begin{teilezwei}
\tz $m + 2n =$ \leerfeld & \tz $m^2 + n =$ \leerfeld \\
\tz $2mn =$ \leerfeld & \tz $m^2 - mn + n^2 =$ \leerfeld \\
\end{teilezwei}
\end{aufgabe}

\begin{aufgabe}{Ich kann Terme mit zwei Variablen zusammenfassen (neu in diesem Jahr). Fasse zusammen. Nur Glieder mit denselben Variablen und denselben Hochzahlen gehören zusammen.}
\beispiel{4x + 3y + 2x &= 6x + 3y}
\begin{gleichungsraster}[1]
\gl{3x + 2y + 4x} & \gl{5m + 4n + 2n} \\ \noalign{\vspace{5mm}}
\gl{2m + 3n + 6m + 4n} & \gl{4u + 5v + 3u + 2v} \\ \noalign{\vspace{5mm}}
\gl{8x - 3y - x + 6y} & \gl{x^2 + 3xy + 2x^2 - xy + x} \\ \noalign{\vspace{5mm}}
\gl{5mn - 2nm + 4m} & \\ \noalign{\vspace{5mm}}
\end{gleichungsraster}
Löse die Klammer auf und fasse zusammen.
\begin{gleichungsraster}[2]
\gl{3 \cdot (2x + y) + 4x} & \\ \noalign{\vspace{5mm}}
\end{gleichungsraster}
\end{aufgabe}

\begin{aufgabe}{Ich kann zwei Klammern miteinander malnehmen. Nimm jedes Glied der ersten Klammer mit jedem Glied der zweiten mal. Schreibe bei a) bis d) nur die vier Produkte auf.}
\beispiel{(x + 6) \cdot (x + 2) &= x^2 + 2x + 6x + 12 \\ &= x^2 + 8x + 12}
\begin{gleichungsraster}[1]
\gl{(x + 1) \cdot (x + 4)} & \gl{(m + 2) \cdot (m + 7)} \\ \noalign{\vspace{5mm}}
\gl{(y + 3) \cdot (y + 1)} & \gl{(u + 4) \cdot (v + 3)} \\ \noalign{\vspace{5mm}}
\end{gleichungsraster}
Multipliziere aus und fasse zusammen.
\begin{gleichungsraster}[2]
\gl{(x + 2) \cdot (x + 9)} & \\ \noalign{\vspace{5mm}}
\end{gleichungsraster}
\end{aufgabe}

\begin{aufgabe}{Ich kann zwei Klammern miteinander malnehmen – weiter. Multipliziere aus und fasse zusammen. Achte auf die Vorzeichen.}
\begin{gleichungsraster}[2]
\gl{(x - 6) \cdot (x + 2)} & \gl{(x - 3) \cdot (x - 7)} \\ \noalign{\vspace{5mm}}
\gl{(2x + 1) \cdot (x + 5)} & \gl{(m + 2n) \cdot (3m + n)} \\ \noalign{\vspace{5mm}}
\gl{(4x - 3) \cdot (2x - 5)} & \\ \noalign{\vspace{5mm}}
\end{gleichungsraster}
\end{aufgabe}

\begin{aufgabe}{Ich finde den Fehler beim Malnehmen zweier Klammern. Finde den Fehler, benenne ihn und korrigiere die Rechnung. Rechne danach die Aufgabe darunter selbst.}
\begin{teile}
\teil Paul rechnet so:
\end{teile}
\rechnung{(x + 6) \cdot (x + 4) &= x^2 + 24}
\begin{gleichungsraster}[2]
\gl{(x + 7) \cdot (x + 2)} & \\ \noalign{\vspace{5mm}}
\end{gleichungsraster}
\begin{teile}
\teil Emma rechnet so:
\end{teile}
\rechnung{(x - 5) \cdot (x - 1) &= x^2 - x - 5x - 5 \\ &= x^2 - 6x - 5}
\begin{gleichungsraster}[2]
\gl{(x - 2) \cdot (x - 8)} & \\ \noalign{\vspace{5mm}}
\end{gleichungsraster}
\begin{teile}
\teil Lukas rechnet so:
\end{teile}
\rechnung{(x + 2) \cdot (x + 4) &= x^2 + 4x + 2x + 8 \\ &= 7x^2 + 8}
\begin{gleichungsraster}[2]
\gl{(y + 6) \cdot (y + 3)} & \\ \noalign{\vspace{5mm}}
\end{gleichungsraster}
\end{aufgabe}

\begin{aufgabe}{Ich kann mit einem Rechteck begründen, warum vier Produkte entstehen. Das Rechteck hat die Seitenlängen $x + 4$ und $x + 1$.}

\flaechenbild{4}{2}{3}{1.2}{x}{4}{x}{1}

\begin{teile}
\teil Schreibe in jedes der vier Teilrechtecke seinen Flächeninhalt.
\teil Begründe mit dem Bild, warum beim Malnehmen von $(x + 4) \cdot (x + 1)$ vier Produkte entstehen.
\teil Tom rechnet $(x + 4) \cdot (x + 1) = x^2 + 4$. Zeige am Bild, welche Teile er vergessen hat.
\end{teile}
\end{aufgabe}

\begin{aufgabe}{Ich kann mit Klammer mal Klammer Flächen vergleichen. Ein quadratisches Beet hat die Seitenlänge $s$~m ($s$ ist größer als 2). Es wird umgebaut: Eine Seite wird $4$~m länger, die andere $2$~m kürzer.}
\begin{teile}
\teil Stelle einen Term für den neuen Flächeninhalt auf und multipliziere aus.
\teil Berechne für $s = 6$ den alten und den neuen Flächeninhalt.
\teil Für welche Seitenlängen wird das Beet durch den Umbau größer? Begründe mit dem Unterschied der beiden Terme.
\end{teile}
\end{aufgabe}
